\documentclass[11pt,a4paper]{article}
\usepackage[T1]{fontenc}
\usepackage[utf8]{inputenc}
\usepackage{lmodern}
\DeclareUnicodeCharacter{1E93}{\d{z}}
\usepackage{amsmath,amssymb}
\usepackage[margin=25mm]{geometry}
\usepackage[hidelinks]{hyperref}
\hypersetup{pdftitle={Newton's force comparison and the renormalization of a trajectory},pdfauthor={navstokgap research working notes}}
\newcommand{\tightlist}{\setlength{\itemsep}{0pt}\setlength{\parskip}{0pt}}
\setlength{\emergencystretch}{3em}
\setcounter{secnumdepth}{0}
\begin{document}
\section{Newton's force comparison and the renormalization of a
trajectory}\label{newtons-force-comparison-and-the-renormalization-of-a-trajectory}

\textbf{6.1 Sol Max, 19 November 2026.} Two ellipses can approach the
same moving circle while their relative displacement still determines
the force gradient in arbitrarily short times. There is a mechanical
requirement that the ellipses can satisfy separately and fail together.
At fixed positive angular momentum, let each have energy \(e>0\) above
the circular minimum. Increase the orbital frequency of the linear
central force from \(\omega\) to \(b\omega\), \(b>1\). Either ellipse
can reach a final ellipse with the same excess \(e\). Under any
\emph{one} change of strength applied to both, at least one finishes
with excess at least \(be\). The exact least possible worst error in
preserving \(e\) is \((b-1)e\), even for smooth positive forces.

This answers two questions about Newton's revision. His comparison of
nascent displacements has a joint classical realization. Implementing an
equal-energy identification by a common drive is a stronger demand with
a sharp classical obstruction. Its absolute cost tends to zero as the
orbits coalesce. For a convergent one-step discretization the classical
evolution is a fixed point of blocking on the space of transfer maps. We
give the scale law for time decimation and insertion of one point, an
exact blocking rule on a physical momentum grid, and a convergent scheme
whose single bare coupling fails to close after one block.

\subsection{What the revision asks of two
motions}\label{what-the-revision-asks-of-two-motions}

Book I, Lemma X of the London \emph{Principia} (1687) specifies a
regular force. In Cambridge (1713) Newton specifies a finite force,
either constant or continually increasing or decreasing. In Corollary 2
he compares departures generated by proportional forces applied at
corresponding parts of similar figures. Corollaries 3 and 4 allow
different forces and times and recover the force ratio from the nascent
displacements divided by the squared times. Proposition VI expressly
gives Corollary 4 as another way to obtain the orbital force formula. In
the preceding Lemma IX he holds the comparison segment \(Ae\) fixed as
the original points coalesce. The segment normalizes the geometry, and
the squared-time ratio in Lemma X compares the intervals. Transporting a
family between strengths at a prescribed energy would be a further
operation.
\href{https://www.newtonproject.ox.ac.uk/view/texts/diplomatic/NATP00077\#par27}{1687,
Lemma X} and
\href{https://www.newtonproject.ox.ac.uk/view/texts/diplomatic/NATP00082\#par27}{1713,
Lemma IX} and
\href{https://www.newtonproject.ox.ac.uk/view/texts/diplomatic/NATP00082\#par31}{1713,
Corollaries 2--4} and
\href{https://www.newtonproject.ox.ac.uk/view/texts/diplomatic/NATP00082\#par90}{Proposition
VI} (diplomatic transcriptions, passages).

For equal masses, the modern vector form uses the departure
\(q(t)-q(0)-t\dot q(0)\) from the inertial tangent. Its coefficient of
\(t^2/2\) is acceleration. Acceleration times the common mass gives
force. Requiring this comparison to work while the two starting states
also approach each other strengthens the historical fixed-pair limit.
The companion formalization derives Corollary 4's ultimate ratio from
the premises of Corollary 3 and a shared proportionality constant. The
formal statement makes explicit where the force law enters.
\href{https://github.com/arivero/newtonlean/blob/8138e9449b4f961c4adbebcb6090acc0d2fd5aaf/NewtonLimitDynamics/Historical/LemmaX/CorollaryIV.lean}{NewtonLean,
Corollary IV at the cited revision} (theorem statements and premises).

\subsection{Two ellipses with a common
clock}\label{two-ellipses-with-a-common-clock}

Use Newton's Proposition X, with acceleration \(-\omega^2q\) in the
plane. Fix the mass \(m>0\), angular frequency \(\omega>0\) and angular
momentum \(J>0\). Choose \(A>B>0\) with \(m\omega AB=J\), and put
\(d=A-B\). The same force acts on the following two test bodies, with
the same origin, axes and time coordinate,

\[
 q_1(t)=(A\cos\omega t,B\sin\omega t),\qquad
 q_2(t)=(B\cos\omega t,A\sin\omega t).
\]

They start at opposite radial turning points on the same ray, with
tangential velocities in the same direction. Their separation has
constant length \(d\), and their velocity difference has length
\(\omega d\). As \(d\) tends to zero they approach the same timed circle
of radius \(\sqrt{J/(m\omega)}\). Both have energy \(E\), circular
minimum \(E_c=\omega J\), and radial action \(I_r=(2\pi)^{-1}\oint
p_r\,dr\). The classical action formula at fixed \(J\) gives

\[
 e=E-E_c=\tfrac12m\omega^2d^2,\qquad
 I_r=\frac{e}{2\omega}=\frac{m\omega d^2}{4}.                       \tag{1}
\]

Here \(I_r\) labels a particular ellipse. The angular action \(J\) stays
fixed in the limit.
\href{https://www.newtonproject.ox.ac.uk/view/texts/diplomatic/NATP00082\#par114}{Newton,
1713, Proposition X} (passage) and
\href{https://arivero.github.io/navstokgap/notes/rivero-1998-conjecture-central-forces.html\#hooke-and-kepler-radial-actions-and-maslov-phases}{classical
radial actions, equations (10)--(11)}.

Write \(\Delta q=q_1-q_2\) and take the component along their common
launch ray \(x\). Lemma X's normalized relative tangent departure is

\[
 \frac{2[\Delta q_x(t)-\Delta q_x(0)-t\Delta\dot q_x(0)]}{d\,t^2}
 =\frac{2(\cos\omega t-1)}{t^2}
 \longrightarrow-\omega^2 .                                        \tag{2}
\]

The error is at most \(\omega^4t^2/12\), independently of \(d>0\). Thus
any joint approach \(d,t\downarrow0\) gives the same force gradient per
unit mass. At \(d=0\) the quotient refers to a limit of pairs. The
mechanical input to its normalization is their actual launch separation.

\textbf{What happens on Newton polygons.} Give both bodies the same
partition \(\pi\) of \([0,t]\). On a cell of length \(h_i\) apply the
impulse before or after free flight, at time \(s_i\) equal to the left
or right endpoint. Recompute both histories when a point is inserted. By
linearity the relative \(x\) coordinate divided by \(d\) is a single
polygon \(u\) starting at \(1\) with zero velocity. Summing the drift
due to each impulse gives

\[
 u(t)=1-\omega^2\sum_i h_i(t-s_i)u(s_i).
\]

The formula with only past kicks also holds at any intermediate time.
Writing \(U=\sup_{[0,t]}|u|\) gives \(U\le1+\omega^2t^2U\). For
\(\omega t\le1/\sqrt2\) this gives \(|u(s)-1|\le2\omega^2t^2\). The
endpoint quadrature for \(\int_0^t(t-s)\,ds\) has error at most
\(\sum_i h_i^2/2\). With \(|\pi|=\max_i h_i\), the two estimates yield

\[
 \left|\frac{2(u(t)-1)}{t^2}+\omega^2\right|
 \le \omega^2\frac{|\pi|}{t}+4\omega^4t^2.                          \tag{3}
\]

It suffices to refine \emph{within} the shrinking interval,
\(|\pi|/t\to0\), with arbitrary \(d>0\). On a single cell of length
\(t\), \(2(u(t)-1)/t^2\) is \(-2\omega^2\) for a kick before flight and
zero for a kick after flight. Both meshes shrink in absolute time. The
general fixed-duration polygon convergence theorem is used with its
common-force Lipschitz hypothesis.
\href{https://arivero.github.io/navstokgap/notes/balaban-ode-lesson.html\#start-with-a-finite-newton-polygon}{Mixed
impulse orders and finite-time convergence}.

\subsection{A requirement that holds separately and fails for the
pair}\label{a-requirement-that-holds-separately-and-fails-for-the-pair}

At the final source frequency \(b\omega\), an ellipse with the same
excess \(e\) would have radial action \(e/(2b\omega)\). Is there one
mechanical change of the source that realizes this comparison for both
launches? Allow any prescribed finite-duration, smooth, positive
isotropic frequency \(\Omega(t)\) joining \(\omega\) to \(b\omega\),
constant near its ends. It is common to both bodies and cannot depend on
which launch it is driving. They exert no force on each other. At
frequency \(\nu\) the Hamiltonian is
\(H_\nu(q,p)=|p|^2/(2m)+m\nu^2|q|^2/2\). Define the final excess for
launch \(i\) by \(e_{f,i}=H_{b\omega}(q_i,p_i)-b\omega J\).

\textbf{Theorem.} For the pair above and \(b>1\),

\[
 \inf_{\Omega}\max_{i=1,2}|e_{f,i}-e|=(b-1)e.                       \tag{4}
\]

One smooth positive drive attains the infimum for all these pairs at
once. For either launch separately, another smooth positive drive can
give \(e_{f,i}=e\).

\textbf{Proof.} At frequency \(\nu\) write \(X=\sqrt{m\nu}\,q\),
\(Y=p/\sqrt{m\nu}\). These coordinates have units of square root of
action and \(H_\nu=\nu(|X|^2+|Y|^2)/2\). The drive has the same
two-by-two matrix \(S\) on \((X_j,Y_j)\) for each Cartesian component
\(j\), using initial and final units at their respective ends. The
Hamiltonian transfer and both changes of units have determinant one. Let
\(l_1,l_2\) be the squared lengths of the columns of \(S\). The two
launches exchange the weights \(m\omega A^2,m\omega B^2\) of those
columns. Consequently their arithmetic mean obeys

\[
 \frac{E_{f,1}+E_{f,2}}2=bE\,\frac{l_1+l_2}{2}\ \ge\ bE,             \tag{5}
\]

since \(l_1+l_2\ge2|\det S|=2\). Subtracting the common final circular
energy gives \(\max_i e_{f,i}\ge be\). This is a comparison of two
specified solutions. The established oscillator inequality for a uniform
ensemble of initial phases uses the same transfer matrix. Equation (5)
identifies two coalescing fixed-\(J\) launches which already realize its
mean exactly. \href{https://arxiv.org/pdf/nlin/0506033}{Robnik and
Romanovski, equations (9)--(12) and the following mean-energy
inequality} (passage).

The lower bound holds even before the trajectory limit. Each free drift
and each isotropic central kick has determinant one on a Cartesian
coordinate and its momentum, and preserves \(J\) in the plane. Any
common partition and common kick placement therefore give (5) exactly
for the two discrete endpoint energies.

For attainment, the positive intermediate frequency
\(\Omega_p=\sqrt b\,\omega\), run for \(\pi/(2\Omega_p)\), has the
quarter-period map \((q,p)\mapsto(p/(m\Omega_p),-m\Omega_p q)\).
Switching to \(b\omega\) gives \(E_f=bE\) for every input. There is also
an exactly smooth realization. Take a smooth positive scale \(\rho\)
constant at 1 initially and at \(b^{-1/2}\) finally. The Ermakov
prescription \(\Omega^2=\omega^2/\rho^4-\ddot\rho/\rho\) preserves the
endpoint oscillator action \(E/\nu\) for every phase. Stretching the
duration makes \(\Omega^2>0\) throughout. Thus both excesses equal
\(be\), which attains (4). \href{https://arxiv.org/pdf/0910.0709}{Chen
et al., scaling equation and endpoint conditions} (passage, used for
compression).

For launch 1 choose instead \(\Omega_p^2=b\omega^2 B/A\). At the end of
its quarter period the position and momentum are perpendicular with
\(|p_f|=mb\omega|q_f|\). The final orbit is the circle and
\(e_{f,1}=0\). For launch 2 interchange \(A,B\). Varying \(\Omega_p\)
continuously from the value for the chosen circle to \(\sqrt b\,\omega\)
therefore brackets the desired excess strictly between 0 and \(be\).
Smooth both switches over a sufficiently short positive time. Continuous
dependence of the linear transfer on the ramps preserves the strict
bracket for this fixed \(e>0\), and continuity in \(\Omega_p\) gives the
target exactly. Time spent at the final strength can equalize the
durations without changing either final energy. This proves the separate
assertion. \(\square\)

At \(b=2\) the best common drive misses the equal-excess target on at
least one orbit by \(e\), one full initial excitation. In final
radial-action units the sharp error is \((1-1/b)I_r\). The fractional
error survives arbitrarily fine comparison of neighbouring trajectories.
The absolute error vanishes quadratically with \(d\). In particular the
limiting moving circle itself admits a common drive to the final circle.

\subsection{A Kepler bound from the whole radial
orbit}\label{a-kepler-bound-from-the-whole-radial-orbit}

For the potential \(-m\mu/r\), take \(m,J,\mu>0\). At fixed \(J\) the
bound radial action \(I\) and energy above the circular minimum obey

\[
 e_\mu(I)=\frac{m^3\mu^2}{2}
       \left(\frac1{J^2}-\frac1{(J+I)^2}\right).
\]

The closed radial curve of excitation \(e\) bounds a disk of canonical
area \(2\pi I\) in \((r,p_r)\). This is the action of the radial
oscillation.
\href{https://arivero.github.io/navstokgap/notes/rivero-1998-conjecture-central-forces.html\#hooke-and-kepler-radial-actions-and-maslov-phases}{Kepler
action and its circular limit} and
\href{https://physique.ensta-paris.fr/PAT/DM/paper_dm3.pdf}{Ramond and
Perez, radial-action formula in the Kepler limit} (passage following
equation (2.7)).

Increase the attraction to \(\lambda\mu\), \(\lambda>1\). For any smooth
rotationally invariant Hamiltonian drive with these endpoints whose
reduced finite-time flow exists on a neighbourhood of the initial closed
disk, at least one initial phase finishes with

\[
                    e_f\ \ge\ \lambda^2 e.                          \tag{6}
\]

The disk argument uses conservation of phase volume, also used for the
classical minimal-work principle. Here the observable is the maximum
over initial phases at fixed radial action.
\href{https://doi.org/10.1103/PhysRevE.75.051124}{Allahverdyan and
Nieuwenhuizen, Phys. Rev.~E 75, 051124} (abstract).

Indeed \(e_{\lambda\mu}(I)=\lambda^2 e_\mu(I)\). The target disk at that
excitation has exactly the initial disk's area. If the entire image of
its boundary lay strictly inside the target disk, its enclosed image
would have strictly smaller area. The reduced Hamiltonian flow preserves
\(dr\,dp_r\). This proves (6), including the case in which a phase
escapes to nonnegative energy. Take \(0<e<m^3\mu^2/(2J^2)\) so that the
initial disk is compact. There are no walls or thermal preparations in
this comparison. The proof finds a phase on the entire radial orbit.
Identifying a prescribed two-phase witness for a nonlinear Kepler drive
would require further information about its transfer.

\subsection{Frequency doubling, decimation and point
insertion}\label{frequency-doubling-decimation-and-point-insertion}

There are two senses in which Newton's construction resembles
renormalization. They admit a direct comparison in this example.

\textbf{Comparing source strengths.} At fixed \(J,e\), the map on the
labels of ellipses is

\[
             (\omega,I_r)\ \longmapsto\ (b\omega,I_r/b).
\]

These maps compose by multiplication of \(b\) and preserve
\(2\omega I_r=e\). For \(b=2\) the radial action halves. This resembles
a Kadanoff change of scale with a specified quantity held fixed. It
compares two different force strengths. Equation (4) gives a mechanical
way to distinguish this map from a common time evolution of the two
orbits. Kadanoff groups spins into cells, and Wilson and Kogut allow
effective interactions to change as degrees of freedom are eliminated.
\href{https://inspirehep.net/files/1b6824b02d6ea885938ba7992c217b69}{Kadanoff,
1966, cell construction} and
\href{https://theory.tifr.res.in/~tridib/ReferenceMaterial/WilsonKogut.pdf}{Wilson
and Kogut, 1974, pp.~82--83 and §12} (passages).

\textbf{Blocking actual trajectory data.} Hold the source and orbit
fixed. At time spacing \(h\), every Cartesian position component
satisfies \(q_{j+1}+q_{j-1}=2c(h)q_j\), where \(c(h)=\cos(\omega h)\).
Keeping even sites gives the exact recurrence at spacing \(2h\),

\[
                   c(2h)=2c(h)^2-1.                                \tag{7}
\]

The decimated data describe the same timed orbit with its original
radial action. Doubling \(\omega\) at fixed \(h\) gives the same
recurrence coefficient by changing the \emph{source}. For instance the
position response to initial momentum in time \(2h\) is
\(\sin(2\omega h)/(m\omega)\), whereas the response at frequency
\(2\omega\) in time \(h\) is \(\sin(2\omega h)/(2m\omega)\). The two
operations agree on the trace of the transfer and can disagree on the
transfer itself.

\textbf{Inserting one point and changing the reference interval.}
Suppose \(0<\omega h<\pi\). Given the two endpoints, the unique
intermediate position at time \(a h\), \(0<a<1\), is

\[
 q(ah)=\frac{\sin((1-a)\omega h)\,q(0)+\sin(a\omega h)\,q(h)}
                 {\sin(\omega h)}.                                  \tag{8}
\]

The subinterval transfers compose to the old endpoint transfer. This is
the classical oscillator instance of the established perfect
discretization construction. The Gell-Mann--Low aspect appears when one
expresses the \emph{same} physics using a coefficient fixed at a new
reference interval. The change of coefficient and its differential form
here are

\[
 \Sigma_s(c)=\cos(s\arccos c),\qquad
 h\,\frac{dc}{dh}=-(\arccos c)\sqrt{1-c^2}.                         \tag{9}
\]

The invariant is \(\arccos(c(h))/h=\omega\). Inserting one point uses
\(\Sigma_a(c)\) and \(\Sigma_{1-a}(c)\) for the two new coefficients.
One has \(\Sigma_s\Sigma_v=\Sigma_{sv}\) wherever the initial,
intermediate and final angles lie in \((0,\pi)\). On that branch
\(\Sigma_2\) is precisely (7). Outside it different frequencies can give
the same sampled coefficient, and at \(\omega h=\pi\) the endpoints do
not determine the inserted point. This supplies an exact
scale-composition example of the Gell-Mann--Low functional equation
idea. The running quantity in (9) is a coefficient of a classical
recurrence. \href{https://doi.org/10.1103/PhysRev.95.1300}{Gell-Mann and
Low, 1954} (abstract on scale functional equations),
\href{https://arxiv.org/pdf/1101.4775}{Bahr, Dittrich and Steinhaus, §3}
(harmonic refinement passage) and
\href{https://arxiv.org/pdf/quant-ph/9803035}{Rivero, 1998, equations
(5)--(10)} (passages on scale control and insertion).

\subsection{How generally does a discretization define a renormalization
flow?}\label{how-generally-does-a-discretization-define-a-renormalization-flow}

\textbf{The continuum evolution as a fixed point.} Let \(A_{t,s}\) be
any continuous one-step rule on the full state for \(s<t\). A blocking
transformation on such \emph{rules} is

\[
 (\mathcal B A)_{t,s}=A_{t,r}\circ A_{r,s},\qquad r=(s+t)/2.
\]

Assume the products on all partitions of each finite interval converge,
as their mesh tends to zero, locally uniformly in the launch state.
Assume the intermediate states stay in compact subsets of the common
domain for compact sets of launches. Write \(\Phi_{t,s}\) for this
limit. Then \((\mathcal B^n A)_{t,s}\to\Phi_{t,s}\), since these are the
products on the dyadic refinements. More strongly, refine \([s,u]\) and
\([u,t]\) separately for any \(s<u<t\). Their union is another
vanishing-mesh partition. Local uniform convergence on the intermediate
compact set and continuity of the limiting transfer give

\[
                  \Phi_{t,s}=\Phi_{t,u}\circ\Phi_{u,s}.
\]

Thus the continuum law is fixed by blocking at \emph{any} inserted
point. When the method converges to Newton's initial-value problem this
fixed law is the classical flow. This argument takes convergence as a
hypothesis. Concrete estimates establishing it for the two Newton
impulse orders are available.
\href{https://arivero.github.io/navstokgap/notes/refinement-composition-and-limit.html\#c.-all-refinement-sequences-and-two-microscopic-kick-orders}{All-partition
construction}.

\textbf{What a change of variables preserves.} Consider any method with
specified observations at each resolution. Let \(\pi'\) refine \(\pi\).
Let \(S_\pi x\) denote what the exact curve \(x\) would give at
resolution \(\pi\), and \(R_{\pi\pi'}\) the restriction of fine data to
coarse data. Require \(R_{\pi\pi'}S_{\pi'}=S_\pi\) and require
successive restrictions to compose. Taking the image of a set of fine
histories under \(R\) then gives exact, associative blocking on sets of
histories. An invertible change of finite data \(C_\pi\) at each
resolution gives the conjugate restriction
\(C_\pi R_{\pi\pi'}C_{\pi'}^{-1}\). It has the same composition law,
including in a full discrete Fourier basis. Retaining selected Fourier
modes introduces a further projection whose closure has to be tested.
Information discarded by restriction can prevent its inversion.

For an approximate method let \(x_\pi\) have error at most \(\eta_\pi\)
relative to \(S_\pi x\) in a chosen norm. If the restriction is
\(L\)-Lipschitz, coarse predictions satisfy

\[
 \|R_{\pi\pi'}x_{\pi'}-x_\pi\|
 \le L\eta_{\pi'}+\eta_\pi .                                       \tag{10}
\]

This triangle inequality gives asymptotic agreement whenever the right
side tends to zero. Sampling a common timed curve in the sup norm has
\(L=1\). Equation (3) illustrates why a normalized observable with
shrinking denominator needs its own estimate. Exact closure on a
proposed set of retained variables \(Pz\) has a further condition. For a
fine block transfer \(F\), \textbf{all fine states with the same \(Pz\)
must have the same \(PFz\)}. This condition is necessary and sufficient
to define the coarse deterministic map on the image of \(P\). Closure in
a prescribed family of couplings also requires the resulting map to
belong to that family.

For a countertest take a convergent kick-then-drift discretization of
the same oscillator. In \((q,\dot q)\) its time-\(h\) step at stiffness
per mass \(k=\omega^2\) is

\[
 M_h(k)=\begin{pmatrix}1-kh^2&h\\-kh&1\end{pmatrix}.
\]

The lower right entry of \(M_h(k)^2\) is \(1-kh^2\). Every
\(M_{2h}(k')\) has lower right entry one. Thus even arbitrary
renormalization of the single stiffness cannot give the blocked
full-state step at finite \(h>0\), \(k>0\). In the space of
determinant-one transfer matrices the exact block is \(M_h(k)^2\). This
example gives both the flow on enlarged data and the failure of a
particular one-coupling description.
\href{https://arivero.github.io/navstokgap/notes/refinement-composition-and-limit.html\#c.-all-refinement-sequences-and-two-microscopic-kick-orders}{Convergence
for arbitrary harmonic partitions}.

\textbf{A grid in physical momentum.} One can use momentum itself as the
mesh coordinate. For one Cartesian harmonic mode of energy \(E_x>0\),
put \(P_*=\sqrt{2mE_x}\). On a monotone arc retain the sign
\(\sigma=\operatorname{sgn}q\). From the energy law
\(q(p)=\sigma\sqrt{P_*^2-p^2}/(m\omega)\). For a cell of momenta
\(p_0,p_1\) traversed in the direction of motion its elapsed time is

\[
 \tau_\sigma(p_0,p_1)=-\frac{\sigma}{\omega}
      \big[\arcsin(p_1/P_*)-\arcsin(p_0/P_*)\big].                  \tag{11}
\]

Inserting any intermediate momentum composes by addition of cell times.
In particular, take midpoint approximations to
\(1/(\omega\sqrt{P_*^2-p^2})\) on a nonuniform momentum grid. Blocking
adjacent cells replaces their inverse speeds in momentum space,
\(dt/|dp|\), by their momentum-width-weighted mean. The rule is
associative and preserves elapsed time. On compact subintervals of
\((-P_*,P_*)\) the midpoint sums converge uniformly for cumulative time.
This extends to the turning points. The integrand is convex and
integrable, so every midpoint sum over a terminal layer is bounded by
the exact integral over that layer enlarged by one mesh cell. This bound
tends to zero with the layer width and mesh. Linear interpolation and
inversion of the strictly monotone cumulative time recover the timed
arc. At a turn one joins signed arcs and keeps their order. Forget the
sign and the same momentum \(p=0\) allows \(q=\pm P_* /(m\omega)\),
giving opposite \(\dot p\). A deterministic rule on that one retained
momentum fails the criterion above. The momentum mesh width has units of
momentum. Neither a time lattice nor a quantum phase was used in
obtaining (11).

\subsection{Consequence for STATE}\label{consequence-for-state}

The bare comparison in Lemma X survives joint refinement of time and of
the two launches. Requiring a single change of force to preserve their
common excess energy would add a condition that already fails in
classical mechanics by (4). Its sharp relative error stays finite as the
circle is approached, while its absolute action cost tends to zero. In
the thermal problem the equilibrium work identity gives the ratio an
entropy consequence and the Nernst heat theorem independently constrains
that entropy. Planck's December 1900 paper introduced the energy element
\(h_P\nu\), with \(\nu\) in cycles per unit time, when counting
distributions of energy among resonators in equilibrium.
\href{https://arivero.github.io/navstokgap/notes/thermal-action-central-force.html\#what-a-quarter-period-change-forces-on-every-newtonian-law}{Thermal
central-force proof} and
\href{https://www.ub.edu/hcub/hfq/sites/default/files/planck_1900_deduccio\%281\%29.pdf}{Planck,
1900, pp.~82--83 in translation} (passage on energy elements).

The mechanical inputs are \(m,J,\omega,d\), the timed source profile and
its duration. The inverse-square comparison adds \(\mu\), of dimension
length cubed per time squared, and its chosen excitation. The time grid
uses a spacing \(h\) in seconds. The momentum grid also fixes the
component energy \(E_x\), hence the momentum amplitude \(P_*\), and uses
cell widths in momentum units. The cut fraction \(a\) and the scale
factors \(b,\lambda\) and \(s\) in (9) are dimensionless. The physical
reduced phase constant \(\hbar\), with \(h_P=2\pi\hbar\), never enters.
Even at fixed dimensionful reference \(J\) the entire construction
admits classical dynamics with zero phase scale. An exclusion from the
trajectory limit would require another principle about which comparisons
the mechanics must realize.
\end{document}
